\chapter{Three Hochschild comparison surfaces: curve cochains, mode
shadows, and critical centres}
\label{chap:three-hochschild-unification}\label{ch:three-hochschild-unification-platonic}

\index{Hochschild cohomology!comparison surfaces|textbf}
\index{chain map!ChirHoch to HH mode}
\index{chain map!HH mode to GF}
\index{Gel'fand--Fuchs cohomology}
\index{Feigin--Frenkel centre!degree-zero comparison}

Three constructions carrying Hochschild language meet a chiral algebra,
but only one is Theorem~H. The local chiral Hochschild cochain complex
on the reference curve is the tangent complex of the chiral MC moduli. The
associative Hochschild complex of a chosen formal-disk mode algebra is
a fibrewise algebraic shadow. The Gel'fand--Fuchs complex is continuous
Lie algebra cohomology of the positive modes; it is not Hochschild
cohomology, and scalar GF cochains appear only after a coefficient
choice.

The comparison maps live only on this formal-disk surface and only to
their images. They do not identify local chiral Hochschild cochains with
the categorical Hochschild cochains of a dg or Calabi--Yau category, with
topological Hochschild homology of a ring spectrum, or with the
product-ambient/Hall Hochschild theories of Vol.~III. Those theories
have different inputs, coefficient objects, and ambient products.
Theorem~H remains the curve-level statement for the derived centre
$Z^{\mathrm{der}}_{\mathrm{ch}}(\cA)=\ChirHoch^\bullet(\cA,\cA)$ on a
fixed reference curve.  Theorem~C remains the Verdier decomposition of
$\mathbf C_g(\cA)=R\Gamma(\overline{\mathcal M}_g,\mathcal Z(\cA))$
after the strict-flat C0 comparison.  A supplied
$\iota_Z^{\mathrm{der}}$, followed by~$H^0$ and C0, connects these
surfaces; Theorem~\ref{thm:scalar-trace-complementarity} records the
further normalized scalar trace.
At critical level the Feigin--Frenkel centre compares the chiral and
mode centres in degree zero; ordinary trivial-coefficient GF
cohomology is not thereby identified with the oper algebra.

\section{Three complexes, three gradings, three operads}
\label{sec:three-three-three}

Fix a chiral algebra~$\cA$ on a smooth projective reference curve~$X$
with vacuum section $\mathbf{1}\in\cA$ and $\cD_X$-module structure.
Choose a formal coordinate at a point of~$X$ and a PBW-filtered
formal-disk mode algebra $A_{\mathrm{mode}}$ associated with that
choice. The non-critical statements assume that $\cA$ lies on the
proved Theorem~H surface: PBW chiral Koszul, $E_\infty$-completed,
finite-type/perfect, and away from critical affine levels and the
non-generic simple-quotient points excluded in
Theorem~\ref{thm:main-koszul-hoch}. Critical-level statements are
confined to~\S\ref{sec:critical-level-ff}. Whenever a scalar
Gel'fand--Fuchs target is used, a continuous invariant functional or a
central-extension coefficient module is part of the data.

\begin{definition}[The three comparison complexes; \ClaimStatusProvedElsewhere]
\label{def:three-hochschild-complexes}
The three complexes are:
\begin{enumerate}[label=\textup{(\roman*)}]
\item \emph{Chiral Hochschild cochains} $\ChirHoch^*(\cA)$: the
 complex of multilinear
 $\cD_{X^{n+2}}$-module maps
 \[
 \ChirHoch^n(\cA) \;=\;
 \mathrm{RHom}_{\cD_{X^{n+2}}}\!\bigl(j_{n!}j_n^*(\cA^{\boxtimes n+2}),\,
 \Delta_*\cA\bigr)
 \]
 on the configuration space $C_{n+2}(X)$, with differential induced
 by collision residues along divisors $D_{ij}=\{z_i=z_j\}$, as in
 Chapter~\ref{chap:deformation-theory}. Spectral data: the formal
 variables $\lambda_1,\ldots,\lambda_{n+1}$ of $\End^{\mathrm{ch}}_\cA$
 attached to relative positions.
\item \emph{Mode-algebra Hochschild cochains}
 $\HH^*(A_{\mathrm{mode}},A_{\mathrm{mode}})$: the classical Hochschild
 complex of the chosen associative formal-disk mode algebra
 $A_{\mathrm{mode}}$. The Zhu algebra and the full mode algebra are
 distinct quotients/completions; a comparison statement must name which
 one is being used. The differential is the standard Hochschild
 differential. The cup product and Gerstenhaber bracket are additional
 structures, not the differential. No curve-level spectral variables
 remain.
\item \emph{Gel'fand--Fuchs continuous cochains}
 $C^*_{\mathrm{cont}}(L_+;M)$: the continuous
 Chevalley--Eilenberg complex of the positive-mode Lie algebra
 $L_+\subset \mathrm{Lie}(A_{\mathrm{mode}})$, computed with respect
 to the topology induced by conformal weight
 \cite{Fuks86,Goncharova73}. The scalar GF complex is the case
 $M=\bC$ with trivial coefficients. Antisymmetrising Hochschild
 cochains naturally gives $A_{\mathrm{mode}}$-valued CE cochains; a
 scalar GF class requires an additional invariant functional or
 coefficient module.
\end{enumerate}
\end{definition}

\begin{remark}[Hochschild variant firewall table; \ClaimStatusProvedHere]
\label{rem:hochschild-variant-firewall-table}
\index{Hochschild cohomology!variant firewall table}
The comparison discipline used in this chapter is the following
three-column firewall.
\begin{center}
\renewcommand{\arraystretch}{1.18}
\begin{tabular}{@{}p{0.27\textwidth}p{0.33\textwidth}p{0.32\textwidth}@{}}
\toprule
Variant & Input category/object & Scope in Vol.~I \\
\midrule
Curve-level chiral Hochschild
$\ChirHoch^\bullet(\cA)$
& Chiral algebra on the reference curve~$X$
& Theorem~H applies only on the PBW, completed, finite-window lane. \\
Chiral coHochschild
$\coHoch^\bullet_{\mathrm{ch}}(C)$
& Chiral coalgebra with diagonal bicomodule
& Transported only after a bimodule--bicomodule equivalence; no arbitrary-coalgebra vanishing. \\
Mode-algebra Hochschild
$\HH^\bullet(A_{\mathrm{mode}})$
& Formal-disk associative mode algebra
& Compared to $\ChirHoch$ by filtered maps onto their images, not by default equivalence. \\
Continuous Lie/Gel'fand--Fuchs cohomology
$H^\bullet_{\mathrm{cont}}(L_+;M)$
& Topological Lie algebra with coefficient module~$M$
& Related only after antisymmetrisation plus coefficient data; scalar GF is not chiral Hochschild. \\
Categorical Hochschild
$\HH^\bullet_{\mathrm{cat}}(\mathcal C)$
& dg/CY category
& Separate Morita/CY theory; not controlled by Theorem~H's curve amplitude. \\
Topological Hochschild
$\mathrm{THH}(R)$
& $E_1$-ring spectrum
& Spectrum-level factorization homology over $S^1$; outside Theorem~H. \\
Hall/product Hochschild
& CoHA, Hall--Drinfeld, or product-ambient geometry
& Vol.~III ambient; a chiral specialization creates a new curve object. \\
\bottomrule
\end{tabular}
\end{center}
Every comparison between two rows requires a named filtered map or
quasi-isomorphism with stated hypotheses. Similar notation is not a
comparison theorem.
\end{remark}

\begin{remark}[Theorem H on the reference curve]
\label{rem:theorem-H-reference-curve-lane}
The symbol $\ChirHoch^\bullet(\cA)$ denotes the curve-level chiral
Hochschild cochain complex of the chiral algebra~$\cA$ on the fixed
reference curve~$X$.  A family datum $H_H(\cA;S)$ supplies a strong
deformation retract and gives
\[
 \operatorname{Supp}H^\bullet(\ChirHoch^\bullet(\cA))\subseteq S.
\]
The Fulton--MacPherson collision calculation, OPE-residue lift,
incidence homotopies, completion, and averaging map are the chain data
which construct~$S$.  Bakalov--De Sole--Kac give bounded Virasoro
support $\{0,2,3\}$ and rank-one even superboson support $\{0,1\}$;
each becomes a curve-chart statement through a bounded-to-chart
quasi-isomorphism.  Thus degree three is an actual Theorem~H column
for the Virasoro family.  Mode-algebra, continuous
Gel'fand--Fuchs, categorical, and topological Hochschild classes enter
through their named comparison maps.
\end{remark}

\begin{remark}[Excluded Hochschild theories]
\label{rem:excluded-hochschild-theories}
\index{Hochschild cohomology!categorical!excluded from Theorem H}
\index{Hall--Drinfeld double!Hochschild scope}
Topological Hochschild homology
$\mathrm{THH}(R)$ is factorization homology of an $E_1$-ring spectrum
over the circle. Its circle action, cyclotomic refinements, and
spectrum-level input are not part of the curve-level
$\ChirHoch^\bullet(\cA)$ complex.

The categorical Hochschild complex
$\HH^\bullet_{\mathrm{cat}}(\mathcal C)$ of a dg category~$\mathcal C$
is the derived endomorphism algebra of the identity functor; with a
compact generator it is Morita equivalent to
$\RHom_{A^e}(A,A)$
(Theorem~\ref{thm:derived-center-hochschild}). For a
Calabi--Yau $d$-category it carries the CY shifted Poisson/BV
structure of the Vol.~III Hochschild calculus. This is not
$\ChirHoch^\bullet(\cA)$ on a curve, and Theorem~H says nothing about
its amplitude.

The product-ambient and Hall Hochschild theories attached to
CoHA/Hall--Drinfeld doubles use product geometry and Hall
convolution. They belong to the Vol.~III ambient. A Stage--2
specialisation functor may produce a chiral algebra on a reference
curve, but the resulting curve-level $\ChirHoch^\bullet$ is a new
chiral object; it is not the categorical or Hall Hochschild complex
itself.
\end{remark}

\begin{remark}[Derived centre and open/closed comparison]
\label{rem:derived-centre-not-physical-bulk-without-comparison}
The chiral derived centre is the cochain-level closed-sector object
\[
 Z^{\mathrm{der}}_{\mathrm{ch}}(\cA)
 := C^\bullet_{\mathrm{ch}}(\cA,\cA)
 = \RHom_{\cA\text{-}\cA}^{\mathrm{fact}}(\cA,\cA),
\]
with its brace dg algebra structure
(Section~\ref{sec:chiral-center-theorem}). It is the endocentral
object through which local closed-sector actions on the boundary
chiral algebra factor. A physical bulk object is not identified with
$Z^{\mathrm{der}}_{\mathrm{ch}}(\cA)$ by notation. Such an
identification is a theorem only after a named open/closed comparison
quasi-isomorphism
\[
 \mathrm{OC}_{\cA}\colon
 \mathcal O^{\mathrm{phys}}_{\mathrm{bulk}}(\cA)
 \xrightarrow{\;\sim\;}
 Z^{\mathrm{der}}_{\mathrm{ch}}(\cA)
\]
has been constructed in the relevant open/closed theory. Without
$\mathrm{OC}_{\cA}$, only the cochain-level derived centre is in
scope.
\end{remark}

\begin{definition}[Presentation-indexed derived centres; \ClaimStatusDefinitional]
\label{def:presentation-indexed-derived-centres}
\index{derived centre!presentation-indexed}
\index{Drinfeld centre!not a chiral derived centre}
The notation \(Z^{\mathrm{der}}\) is always presentation-indexed.
For the standard surfaces used in this volume we write:
\begin{enumerate}[label=\textup{(\roman*)}]
\item the curve-level factorisation derived centre
\[
 Z^{\mathrm{der}}_{\mathrm{ch}}(\cA)
 :=
 \RHom_{\cA\text{-}\cA}^{\mathrm{act},\mathrm{fact}}(\cA,\cA)
 \simeq C^\bullet_{\mathrm{ch}}(\cA,\cA),
\]
the brace dg algebra of chiral Hochschild cochains;
\item the ordered \(\Eone\)-chiral derived centre
\[
 Z^{\mathrm{der}}_{\Eone\text{-}\mathrm{ch}}(\cA^{\mathrm{ord}})
 :=
 \RHom_{\cA^{\mathrm{ord}}\text{-}\cA^{\mathrm{ord}}}
 ^{\mathrm{act},\mathrm{ord}}(\cA^{\mathrm{ord}},
 \cA^{\mathrm{ord}}),
\]
computed in the ordered factorisation-bimodule ambient on
\(\Conf_n(X)\);
\item the symmetric or vertex derived centre
\[
 Z^{\mathrm{der}}_{\Einf\text{-}\mathrm{ch}}(\cA^{\mathrm{sym}})
 :=
 \RHom_{\cA^{\mathrm{sym}}\text{-}\cA^{\mathrm{sym}}}
 ^{\mathrm{act},\mathrm{sym}}(\cA^{\mathrm{sym}},
 \cA^{\mathrm{sym}}),
\]
computed after the symmetric chiral product, or after the
\(R\)-twisted ordered-to-symmetric descent when that descent has
actually been constructed;
\item the formal-disk mode derived centre
\[
 Z^{\mathrm{der}}_{\mathrm{mode}}(A_{\mathrm{mode}})
 :=
 \RHom_{A_{\mathrm{mode}}^e}
 (A_{\mathrm{mode}},A_{\mathrm{mode}}),
\]
the associative Hochschild cochain complex of the chosen mode algebra;
\item the categorical Hochschild centre
\[
 Z^{\mathrm{der}}_{\mathrm{cat}}(\mathcal C)
 :=
 \RHom_{\mathrm{Fun}(\mathcal C,\mathcal C)}
 (\mathrm{id}_{\mathcal C},\mathrm{id}_{\mathcal C}).
\]
The braided Drinfeld centre
\(\mathfrak Z_{\mathrm{Dr}}(\mathcal C)\) is a categorical object,
not a chiral derived-centre complex; its Hochschild cochains require
a separate categorical comparison.
\end{enumerate}
These five objects are not identified by notation.  A comparison
between any two of them must name its presentation, ambient category,
completion, and coefficient object.
\end{definition}

\begin{construction}[Presentation-centre comparison maps;
\ClaimStatusConditional]
\label{constr:presentation-centre-comparison-maps}
\index{derived centre!comparison maps}
The comparison maps used in this chapter are the following typed maps,
defined only after the indicated auxiliary data are fixed:
\begin{align}
\rho_{\mathrm{form}}\colon
 Z^{\mathrm{der}}_{\mathrm{ch}}(\cA)
 &\longrightarrow
 Z^{\mathrm{der}}_{\mathrm{mode}}(A_{\mathrm{mode}}),
\label{eq:presentation-centre-formal-map}\\
\operatorname{av}_{R}^{Z}\colon
 Z^{\mathrm{der}}_{\Eone\text{-}\mathrm{ch}}(\cA^{\mathrm{ord}})
 &\longrightarrow
 Z^{\mathrm{der}}_{\Einf\text{-}\mathrm{ch}}(\cA^{\mathrm{sym}}),
\label{eq:presentation-centre-averaging-map}\\
\operatorname{Mor}_{b}^{Z}\colon
 Z^{\mathrm{der}}_{\mathrm{ch}}(A_b)
 &\longrightarrow
 Z^{\mathrm{der}}_{\mathrm{cat}}(\mathcal C_{\mathrm{op}}).
\label{eq:presentation-centre-morita-map}
\end{align}
Here \(\rho_{\mathrm{form}}\) is formal-disk restriction followed by
the PBW filtered mode model and quotient by the spectral-variable
kernel; \(\operatorname{av}_{R}^{Z}\) is the ordered-to-symmetric
centre map induced by the full \(R\)-twisted descent/Reynolds-kernel
package; and \(\operatorname{Mor}_{b}^{Z}\) is the compact-generator
Morita comparison after choosing \(b\in\mathcal C_{\mathrm{op}}\) and
checking completion compatibility.

Each map is a quasi-isomorphism only on its named generic comparison
surface:
\begin{description}[leftmargin=2.0em, style=nextline]
\item[H1.]
PBW chiral Koszulity, an \(E_\infty\)-completed model, finite
type/perfectness, and exclusion of critical or non-generic quotient
points for the formal-disk comparison.
\item[H2.]
Strict flat or Mittag--Leffler completion together with the Theorem~H
amplitude/concentration package, so that the centre complex is being
computed in the same completed ambient on both sides.
\item[H3.]
Full \(R\)-twisted ordered-to-symmetric descent, including the
Reynolds-kernel ideal and compatibility with brace operations.
\item[H4.]
Compact-generator Morita equivalence
\(\mathcal C_{\mathrm{op}}\simeq\Perf(A_b)\), categorical Hochschild
comparison, and compatibility of all completions.
\end{description}
When the relevant hypotheses fail, the displayed arrows remain at
most filtered maps to images or quotients.  In particular, the
factorisation derived centre, ordered centre, symmetric centre, mode
centre, and categorical Hochschild centre do not coincide by default;
the chiral Deligne--Tamarkin rectification changes the operadic
structure on the brace complex and is not itself an equality of
presentation-indexed centres.
\end{construction}

\begin{remark}[Five-object firewall specialization]
\label{rem:five-object-firewall-three-hochschild}
By Convention~\ref{conv:master-five-object-firewall} and
Theorem~\ref{thm:typed-verdier-koszul-firewall}, the five objects
\[
 \cA,\qquad
 B^{\mathrm{ch}}(\cA)=T^c(s^{-1}\bar{\cA}),\qquad
 \cA^{\mathrm i}=C_X(s^{-1}V,s^{-2}R),\qquad
 \cA^!=\mathbb D_X(\cA^{\mathrm i}),\qquad
 Z^{\mathrm{der}}_{\mathrm{ch}}(\cA)
\]
have distinct types.  The chosen presentation carries
$q_\cA\colon\cA^{\mathrm i}\to B^{\mathrm{ch}}(\cA)$, whose
quasi-isomorphism is quadratic Koszul recognition. Bar--cobar inversion is the statement that
$\Omega^{\mathrm{ch}}B^{\mathrm{ch}}(\cA)$ recovers~$\cA$ on the
inversion surface. Verdier Koszul duality produces~$\cA^!$ from the
quadratic coalgebra~$\cA^{\mathrm i}$. The derived centre is the
Hochschild cochain object.  The displayed maps record every passage
between these carriers.
\end{remark}

\begin{remark}[Distinguishing gradings]
\label{rem:three-gradings}
The three complexes carry \emph{different} gradings:
\begin{itemize}[nosep]
\item $\ChirHoch^*$ has cochain degree $n$ and an independent
 \emph{spectral filtration} by collision depth
 $p = |\{i<j : z_i = z_j\}|$;
\item $\HH^*(A_{\mathrm{mode}})$ has cochain degree $n$ and a
 conformal-weight grading inherited from $A_{\mathrm{mode}}$;
\item $H^*_{\mathrm{cont}}(L_+;M)$ has antisymmetric
 Lie-cochain degree $n$ and an outer-derivation filtration.
\end{itemize}
Chain-level comparison maps must respect the cochain degree and are
filtered for the secondary gradings. They are not inverse equivalences
of complexes. The maps below land in associated images; the scalar
GF versions require the coefficient datum named above.
\end{remark}

\section{Chain maps: construction}
\label{sec:three-chain-maps}

\begin{construction}[Filtered map $\Theta_1$:
ChirHoch $\to$ HH$_{\mathrm{mode}}$; \ClaimStatusConditional]
\label{constr:theta1-chirhoch-to-HHmode}
\index{Theta 1 intertwiner@$\Theta_1$ intertwiner}
On the Koszul locus at non-critical level, and after the formal
coordinate and PBW-filtered mode model have been fixed, restrict a
chiral Hochschild cochain $\phi\in\ChirHoch^n(\cA)$ to the formal
disk, extract the constant term in the spectral variables
$\lambda_i$, and project to the chosen mode algebra. Explicitly, on
the associated graded of the PBW filtration:
\[
 \Theta_1(\phi)(a_1\otimes\cdots\otimes a_n)
 \;:=\;
 \mathrm{ev}_0\!\bigl[\phi(\widetilde a_1(z_1),\ldots,
 \widetilde a_n(z_n),\mathbf{1}(z_{n+1}))\bigr]\!
 \Bigr|_{\lambda_i=0}
\]
where $\widetilde a_i\in\cA$ is any formal-disk lift of
$a_i\in A_{\mathrm{mode}}$ and $\mathrm{ev}_0$ evaluates at the
colliding configuration $z_1=\cdots=z_{n+1}$. The construction is a
chain map only under the filtered PBW comparison identifying the
formal-disk restriction of the chiral bar differential with the
Hochschild differential of $A_{\mathrm{mode}}$ after quotienting by
the spectral-variable kernel. Without this comparison, it is only a
filtered projection of cochains.
\end{construction}

\begin{construction}[Antisymmetrisation $\Theta_2$:
HH$_{\mathrm{mode}} \to$ CE; \ClaimStatusConditional]
\label{constr:theta2-HHmode-to-GF}
\index{Theta 2 intertwiner@$\Theta_2$ intertwiner}
The associative mode algebra $A_{\mathrm{mode}}$ generates a Lie
algebra $\mathrm{Lie}(A_{\mathrm{mode}})$ via the commutator
bracket. The chosen mode functor sends the positive modes of~$\cA$ to a
closed Lie subalgebra
$L_+\subset \mathrm{Lie}(A_{\mathrm{mode}})$, equipped with the
conformal-weight topology. If the PBW mode model is faithful on
positive modes, this subalgebra identifies with
$\mathrm{Lie}(\cA)_{>0}$; otherwise $L_+$ is the image after quotienting
by the positive-mode kernel. The antisymmetrization
map
\[
 \mathrm{Alt}\colon
 C^n_{\mathrm{Hoch}}(A_{\mathrm{mode}},A_{\mathrm{mode}})
 \;\longrightarrow\;
 C^n_{\mathrm{cont}}(L_+;A_{\mathrm{mode}}^{\mathrm{ad}})
\]
defined by $\mathrm{Alt}(\phi)(x_1\wedge\cdots\wedge x_n)
= \sum_{\sigma\in S_n}\mathrm{sgn}(\sigma)\,
\phi(x_{\sigma(1)},\ldots,x_{\sigma(n)})$ is a chain map of
complexes from Hochschild cochains to Chevalley--Eilenberg cochains
with adjoint coefficients. Given a continuous invariant functional
$\ell\colon A_{\mathrm{mode}}\to M$, or the critical oper coefficient
module in~\S\ref{sec:critical-level-ff}, composition gives a scalar or
$M$-valued map
\[
 \Theta_{2,\ell}:=\ell\circ\mathrm{Alt}\colon
 C^n_{\mathrm{Hoch}}(A_{\mathrm{mode}},A_{\mathrm{mode}})
 \longrightarrow C^n_{\mathrm{cont}}(L_+;M).
\]
There is no canonical map to trivial-coefficient scalar GF cochains
without this coefficient datum. Continuity of the image follows from
the assumption that $\phi$ is polynomial in its mode arguments.
\end{construction}

\begin{construction}[Composite $\Theta_{3,\ell}$:
ChirHoch $\to$ CE/GF image; \ClaimStatusConditional]
\label{constr:theta3-chirhoch-to-GF}
\index{Theta 3 intertwiner@$\Theta_3$ intertwiner}
After choosing the PBW comparison and coefficient datum~$\ell$, define
$\Theta_{3,\ell}:=\Theta_{2,\ell}\circ\Theta_1$. Equivalently,
evaluate a chiral Hochschild cochain on the formal-disk restriction,
restrict further to the positive-mode subalgebra~$L_+$, extract the
constant term in the spectral variables, antisymmetrize, and apply the
coefficient functional. The composition is a chain map precisely under
the hypotheses making both factors chain maps.
\end{construction}

\begin{remark}[No chain map back; asymmetry is structural]
\label{rem:no-chain-map-back}
There is no canonical chain map in the opposite direction. The
positive-mode subalgebra $L_+$ loses the mode $a_{(0)}$ content
encoded in chiral derivations at simple pole order; the Zhu
functor loses the spectral parameters; lifting $\HH^*$ cochains
to $\ChirHoch^*$ requires choosing formal primitives that are
not canonical. The cohomological statements below are image
statements after quotient and coefficient projection, not inverses on
the full truncated complexes.
\end{remark}

\section{Chain-level image comparison in low degrees}
\label{sec:chain-level-agreement-low}

\begin{theorem}[Image-level comparison in low degrees on the Koszul
locus; \ClaimStatusConditional]
\label{thm:three-hochschild-chain-level-agreement-low-degree}
\index{Theorem!three-Hochschild chain-level image comparison}
Let $\cA$ lie on the proved Theorem~H surface named above. Assume the
formal-disk PBW comparison makes $\Theta_1$ a filtered chain map after
quotienting by the spectral-variable kernel, assume the positive-mode
image~$L_+$ and coefficient datum~$\ell$ are fixed for the CE/GF target,
and let $\tau_{\le 2}$ denote stupid truncation to cochain degrees
$\{0,1,2\}$. Then the chosen maps compare the low-degree surfaces only
after passing to the indicated images/quotients:
\begin{enumerate}[label=\textup{(\roman*)}]
\item $\Theta_1$ induces a quasi-isomorphism from the quotient of
 $\tau_{\le 2}\ChirHoch^*(\cA)$ by the spectral-variable kernel onto
 its image in $\tau_{\le 2}\HH^*(A_{\mathrm{mode}})$;
\item $\Theta_{2,\ell}\colon
 \tau_{\le 2}\HH^*(A_{\mathrm{mode}}) \to
 \tau_{\le 2} C^*_{\mathrm{cont}}(L_+;M)$ is a chain map whose
 antisymmetric quotient is quasi-isomorphic to its continuous-cochain
 image. When $M=\bC$ with trivial coefficients, this image is only the
 part of scalar GF cohomology detected by~$\ell$;
\item $\Theta_{3,\ell} = \Theta_{2,\ell}\circ\Theta_1$ is therefore a
 chain-level comparison from the chiral quotient to the selected
 continuous-cochain image. It is not asserted to be a quasi-isomorphism
 from the full truncated chiral Hochschild complex, nor onto full
 scalar GF cohomology.
\end{enumerate}
\end{theorem}

\begin{proof}
For~(i), quotienting by the spectral-variable kernel is exactly the
operation that removes the part of the chiral cochain invisible to the
formal-disk mode algebra. The filtered PBW comparison hypothesis then
identifies the induced differential on this quotient with the
Hochschild differential on the image. Hence the quotient maps
quasi-isomorphically onto that image. Theorem~H supplies only the
curve-level cohomological amplitude; it is not used to identify the
mode algebra or to kill the kernel.

For~(ii), antisymmetrisation is the standard Hochschild-to-Chevalley
map from an associative algebra to its commutator Lie algebra with
adjoint coefficients. Composing with~$\ell$ remains a chain map exactly
when~$\ell$ is invariant, or more generally when the chosen coefficient
module makes it a morphism of CE complexes. Thus the statement is an
image statement. It does not assert that every scalar GF class in
degrees~$\le 2$ comes from Hochschild cochains.

For~(iii), composing the quotient comparison in~(i) with the
antisymmetric-image comparison in~(ii) gives the asserted map from the
chiral quotient to the selected continuous-cochain image. The explicit
tables below show why no statement about the full unquotiented
complexes can be true.
\end{proof}

\begin{remark}[Explicit low-degree tables; \ClaimStatusConditional]
\label{rem:three-hochschild-low-degree-tables}
The two bounded computations with primary-source support are as
follows.  The last column is a curve-chart conclusion precisely when
the displayed comparison is a quasi-isomorphism.
\begin{center}
\renewcommand{\arraystretch}{1.15}
\begin{tabular}{lll}
\toprule
family & bounded vertex cohomology & chart transport \\
\midrule
$B_{\mathfrak h}$, $\dim\mathfrak h_{\bar0}=1$ &
$\dim H^n_{\mathrm{ch},b}=(2,1,0,\ldots)$ &
$\chi_{B_{\mathfrak h}}^{\mathrm{bd}}$ \\
$\operatorname{Vir}_c$ &
$\dim H^n_{\mathrm{ch},b}=1$ for $n\in\{0,2,3\}$ &
$\chi_{\operatorname{Vir}_c}^{\mathrm{bd}}$ \\
\bottomrule
\end{tabular}
\end{center}
These are \cite[Theorems~7.4 and~7.2]{BDSK21}.  The Virasoro row
contains a genuine degree-three class.  Ordinary mode-algebra
Hochschild cohomology and continuous Gel'fand--Fuchs cohomology enter
through $\Theta_1$, $\Theta_{2,\ell}$, and $\Theta_{3,\ell}$; their
dimension vectors become comparable only on the quasi-isomorphism
range of those maps.  For a chosen quadratic presentation, the
bar-dual generator is the degree-one part of
$\cA^i=C_X(s^{-1}V,s^{-2}R)$, carried to the full bar by~$q_\cA$.
\end{remark}

\begin{remark}[What ``agreement'' means precisely; \ClaimStatusConditional]
\label{rem:three-hochschild-agreement-precision}
Type signature: \textup{(}Open quadrant, low-degree Hochschild
crosswalk comparison, Beilinson level~$3$ derived-centre presentation,
hypothesis package: truncated PBW/coefficient comparison plus the
separate spectrum-level lift package when THH is invoked\textup{)}.
Three distinct comparison strengths occur:
\begin{enumerate}[label=\textup{(\alph*)}]
\item \emph{Cohomological dimensions.}  These are invariants of the
native complexes and agree only after a comparison quasi-isomorphism
has been constructed.
\item \emph{Low-degree Euler characteristics.}  These depend on the
coefficient convention; trivial-coefficient GF, adjoint CE, and a
scalar anomaly coefficient give different traces.
\item \emph{Chain-level quasi-isomorphism onto image via
 $\Theta_1,\Theta_{2,\ell},\Theta_{3,\ell}$.} Conditional on the PBW
 and coefficient hypotheses above, and only on the truncated range
 $\tau_{\le 2}$; this is
 Theorem~\ref{thm:three-hochschild-chain-level-agreement-low-degree}.
\end{enumerate}
Statement~(c) is the chain-level formulation used here.  Statements
(a) and~(b) become consequences only on its quasi-isomorphism range
and with the coefficient convention fixed. This is a structural
distinction among operads and differentials. A THH comparison enters
after the spectrum-level lift
package of Remark~\ref{rem:hochschild-crosswalk-spectrum-lift-package}
has been supplied; the high-degree comparison target used below is
scalar GF.
\end{remark}

\section{Cohomological image comparison on a specified range}
\label{sec:three-hochschild-all-degree}

\begin{theorem}[Family-supported cohomological image comparison;
\ClaimStatusConditional]
\label{thm:three-hochschild-cohomological-agreement-all-degree}
\index{Theorem!three-Hochschild cohomological image comparison}
Let $\cA$ carry $H_H(\cA;S)$, and let $J\subset\mathbb Z$ be a set of
degrees.  Assume that the filtered maps
$\Theta_1$ and $\Theta_{2,\ell}$ of
Theorem~\ref{thm:three-hochschild-chain-level-agreement-low-degree}
are quasi-isomorphisms from their quotient complexes onto their images
after truncation to~$J$.  Then the induced maps
\[
 [\Theta_1]\colon H^n\ChirHoch^*(\cA)
 \to H^n\HH^*(A_{\mathrm{mode}}),
 \qquad
 [\Theta_{2,\ell}]\colon H^n\HH^*(A_{\mathrm{mode}})
 \to H^n_{\mathrm{cont}}(L_+;M),
\]
are isomorphisms from the corresponding quotient cohomologies onto
their respective images for $n\in J$. The image of
$[\Theta_1]$ in $H^n\HH^*$ is the mode-visible part of the chiral
class. The image of $[\Theta_{2,\ell}]$ in continuous cohomology is
the $\ell$-visible part.  The family retract also gives
$H^n(\ChirHoch^\bullet(\cA))=0$ for $n\notin S$.  If $S\subseteq J$,
the comparison range accounts for every nonzero chiral source group.
\end{theorem}

\begin{proof}
A quasi-isomorphism of truncated quotient complexes induces an
isomorphism on their cohomology in the truncation range.  Apply this
first to $\Theta_1$ and then to the coefficient-selected
antisymmetrization $\Theta_{2,\ell}$.  The support assertion is
Theorem~H for $H_H(\cA;S)$.
\end{proof}

\begin{remark}[Family support and Gel'fand--Fuchs support]
\label{rem:GF-unbounded-but-image-bounded}
The continuous GF cohomology $H^*_{\mathrm{cont}}(L_+;\bC)$
is a polynomial ring for Virasoro (\cite[Theorem~1.4.3]{Fuks86}:
$H^*(L_+;\bC) = \bC[c_2]$ with $|c_2|=2$;~\cite{Goncharova73}) and
unbounded in general. This is the support of the mode Lie algebra
complex. The image of
$\Theta_{3,\ell} = \Theta_{2,\ell}\circ\Theta_1$ is the
$\ell$-visible part in the range where the
antisymmetrization of $\HH^*_{\mathrm{mode}}$ sees continuous-Lie
content after the coefficient choice.  If $\cA$ carries
$H_H(\cA;S)$, the chiral source occurs only in~$S$, whereas the GF
target retains its native support. The chosen mode algebra determines
$H^n\HH^*(A_{\mathrm{mode}})$.
\end{remark}

\section{Support separation beyond the family set}
\label{sec:three-hochschild-divergence}

\begin{proposition}[Comparison cone beyond the family support;
\ClaimStatusConditional]
\label{prop:three-hochschild-high-degree-divergence}
\index{divergence!three Hochschild at n>=3}
Let $\cA$ carry $H_H(\cA;S)$, and let
$\Theta_{3,\ell}\colon C^\bullet_{\mathrm{ch}}(\cA,\cA)	o
C^\bullet_{\mathrm{cont}}(L_+;M_\ell)$ be the coefficient-selected
comparison.  Suppose $n,n+1\notin S$ and
$H^n_{\mathrm{cont}}(L_+;M_\ell)\neq0$.  Then
\[
 H^n\!\left(\operatorname{Cone}(\Theta_{3,\ell})\right)
 \cong H^n_{\mathrm{cont}}(L_+;M_\ell)\neq0.
\]
Thus the comparison cone records every target class lying beyond the
family support.  For the Virasoro bounded-to-chart datum
$S=\{0,2,3\}$, the unbounded Goncharova--Fuks continuous cohomology
produces such classes in arbitrarily high degree
\cite{Goncharova73,Fuks86}.
\end{proposition}

\begin{proof}
Theorem~H gives zero source cohomology in degrees $n$ and $n+1$.
The long exact sequence of the displayed cone therefore identifies
its degree-$n$ cohomology with the target group.  The final assertion
uses the unbounded continuous cohomology calculation of
Goncharova--Fuks.
\end{proof}

\begin{remark}[The two native filtrations]
\label{rem:why-theorem-H-not-GF}
Theorem~H uses a family retract of the Fulton--MacPherson
collision-depth filtration on the curve.  Gel'fand--Fuchs cohomology
uses the Chevalley--Eilenberg filtration of a positive-mode Lie
algebra with a chosen coefficient module.  The comparison map retains
the intersection of these two chain structures; its cone measures the
remaining target classes.  Topological, categorical, and
product-ambient Hall Hochschild theories have their corresponding
native filtrations and enter through their own comparison maps.
\end{remark}

\section{Critical-level centres via Feigin--Frenkel}
\label{sec:critical-level-ff}

\begin{theorem}[Critical-level degree-zero comparison;
\ClaimStatusConditional]
\label{thm:critical-level-ff-center-unification}
\index{Feigin--Frenkel centre!degree-zero centre comparison}
\index{critical level!Hochschild centre comparison}
At the critical level $k = -h^\vee$ of the affine vertex algebra
$V_{-h^\vee}(\fg)$, the chiral centre and the associative mode centre
identify with the Feigin--Frenkel centre:
\[
 \ChirHoch^0(V_{-h^\vee}(\fg))
 \;\cong\;
 \HH^0(A_{\mathrm{mode}})
 \;\cong\;
 \mathrm{Fun}(\mathrm{Op}_{\fg^\vee}(D))
\]
where the right-hand side is the Feigin--Frenkel centre, the polynomial
ring of $\fg^\vee$-opers on the formal disk. Ordinary continuous
Gel'fand--Fuks degree-zero cohomology with trivial coefficients remains
$\mathbb C$; it matches the oper algebra only after replacing the
coefficient system by the critical central extension/oper coefficient
module. That coefficient replacement is the conditional part of the
comparison.
\end{theorem}

\begin{proof}
At the critical level the Sugawara denominator degenerates and the
generic comparison acquires the curvature-residue class
$\mathfrak c_{\mathrm{crit}}(\fg)$.  The degree-zero groups are
centres of their native algebras:
\begin{itemize}[nosep]
\item $\ChirHoch^0(V_{-h^\vee}(\fg)) = Z_{\mathrm{ch}}(V_{-h^\vee}(\fg))$
 equals the chiral centre, which Feigin--Frenkel
 \cite{FeFr92,Frenkel07} identify with
 $\mathrm{Fun}(\mathrm{Op}_{\fg^\vee}(D))$.
\item $\HH^0(A_{\mathrm{mode}}) = Z(A_{\mathrm{mode}})$ is the
 classical associative centre of the mode algebra; under the
 Feigin--Frenkel isomorphism the Segal--Sugawara generators of the
 mode centre match the oper-polynomial generators.
\item Ordinary $H^0_{\mathrm{cont}}(L_+;\bC)$ with trivial
 coefficients is $\bC$. Continuous cochains with coefficients in the
 critical central extension/oper module carry the oper algebra; the
 required splitting is the conditional coefficient-change map.
\end{itemize}
The first two identifications share the same Feigin--Frenkel target and
the same generators. The stated coefficient replacement gives the
third comparison.
\end{proof}

\begin{remark}[Scope of the centre comparison]
\label{rem:ff-unification-degree-zero-only}
In higher degrees each comparison surface retains its native
differential and operations. The classical $\HH^*$ of
$A_{\mathrm{mode}}$ carries its Poisson structure; scalar GF
cohomology is a trivial-coefficient CE theory; CY categorical and Hall
Hochschild objects have their product ambients. The Feigin--Frenkel
statement identifies their named degree-zero centre targets.
\end{remark}

\begin{remark}[Chain vs cohomological boundary]
\label{rem:three-hochschild-chain-cohomology-boundary}
The boundary has three strata:
\begin{itemize}[nosep]
\item $n=0$: at critical level, the chiral centre and mode centre
 share the Feigin--Frenkel oper algebra; ordinary
 trivial-coefficient GF degree zero remains~$\bC$.
\item $n=1,2$ at non-critical level: chain-level quasi-isomorphism
 onto image via $\Theta_1,\Theta_{2,\ell},\Theta_{3,\ell}$ under the
 PBW and coefficient hypotheses.
\item for a family datum $H_H(\cA;S)$, the chiral source is supported
 in~$S$; scalar GF and mode Hochschild retain their respective
 supports.  Proposition~\ref{prop:three-hochschild-high-degree-divergence}
 computes the cone wherever the target continues beyond~$S$.
\end{itemize}
This is the sharp statement.
\end{remark}

\section{Formal-disk and cross-ambient boundary}
\label{sec:three-hochschild-consequences}

\begin{corollary}[Formal-disk comparison with fixed coefficients]
\label{cor:three-hochschild-scoped-comparison}
\ClaimStatusConditional
On the non-critical Koszul locus, the formal-disk comparison gives
chain maps
$\Theta_1,\Theta_{2,\ell},\Theta_{3,\ell}$
(Constructions~\ref{constr:theta1-chirhoch-to-HHmode},
\ref{constr:theta2-HHmode-to-GF},
\ref{constr:theta3-chirhoch-to-GF}) only after the PBW comparison,
quotient by the spectral-variable kernel, and coefficient datum~$\ell$
have been fixed. On $\tau_{\le 2}$ these maps are
quasi-isomorphisms onto their images.  For a family datum
$H_H(\cA;S)$, Theorem~H supplies support in~$S$, and
Proposition~\ref{prop:three-hochschild-high-degree-divergence}
computes the comparison cone beyond that set.

At critical level $k=-h^\vee$, the Feigin--Frenkel centre compares
the chiral centre and the mode centre in degree zero
(Theorem~\ref{thm:critical-level-ff-center-unification}); ordinary
trivial-coefficient GF degree zero is~$\bC$; the critical oper
coefficient system gives the stated coefficient-change comparison.
Categorical Hochschild cochains and product-ambient Hall Hochschild
theories enter through their own comparison functors.
\end{corollary}

\begin{remark}[Cross-ambient boundary]
\label{rem:three-hochschild-propagation}
Volume~II may interpret the three formal-disk maps as boundary
truncation, commutator-Lie descent, and anomaly projection. That
interpretation is physical language for the same image statements, not
a new equivalence of Hochschild theories. Vol.~III uses categorical
Hochschild complexes of CY categories and Hall--Drinfeld double
constructions in their own ambient. A Stage--2 chiral specialisation
can produce a curve algebra to which Theorem~H applies, but the
categorical/Hall Hochschild object is not thereby a Vol.~I
consequence.
\end{remark}
